\chapter{Syntax: Modular Agglutinative Structuralism}

\section{Syntactic Invariants: The SOV Constraint}
The syntactic architecture of the PED reconstruction is defined by a rigid Subject-Object-Verb (SOV) template. This structural invariant provided a stable framework for the transmission of linguistic information over deep-time horizons. In this system, the agential Subject (Active) and the resultative Object (Inactive) maintained a fixed relational geometry, which mitigated the impact of phonetic drift on semantic clarity.

\section{The Agglutinative Modeling of Meaning}
Grammatical complexity in PED was achieved through agglutinative stacking, where discrete semantic and functional units were appended to the root in a deterministic sequence. This modular approach allowed for the construction of complex propositions through a series of additive logical layers.

\subsection{Standard Suffix Linearization}
The sequence of morphological components followed a consistent right-branching logic:
\begin{center}
\textbf{[ROOT] + [ANIMACY/ASPECT] + [CASE/TARGET] + [PERSON/RELATION]}
\end{center}

\section{Logical Connectivity: The Converb Analysis}
In the absence of subordinating conjunctions, complex clausal relationships were mediated through converbal structures. These were formed by applying nominal case markers directly to verbal stems to indicate temporal or logical dependency.
\begin{itemize}
    \item \textbf{Temporal Invariant (*-R):} The locative marker applied to a verb stem denoted simultaneous or prior action (e.g., "While [Action]").
    \item \textbf{Directional Invariant (*-K):} The dative marker applied to a verb stem denoted purposive intent (e.g., "In order to [Action]").
    \item \textbf{Conditional Invariant (*-i):} Denoted hypothetical or contingent states (e.g., "Assuming [Action]").
\end{itemize}

\section{Structural Glossing of Reconstructed Units}
We provide systematic glossing of PED syntactic units to illustrate the modularity of the ancestral system, utilizing the empirically derived phonological invariants.

\subsection{Basic Agential Clause}
\begin{center}
\textit{*H-eK tat-N bhar-S-mi} \\
\small [1.SG.ACT THAT-INACT CARRY-PFV-1.SG] \\
\normalsize "I carried that."
\end{center}

\subsection{Spatial Deictic Clause}
\begin{center}
\textit{*H-tV id-N K-N-S-ti} \\
\small [2.SG.ACT THIS-INACT SEE-PFV-3.SG] \\
\normalsize "Thou sawest this."
\end{center}

\section{Phonetic Boundary Rules (Sandhi)}
The intersection of modular suffixes produced predictable phonetic modifications at morpheme boundaries. The primary rule identified is the \textbf{Nasal Assimilation Rule}:
\begin{center}
\textit{*N + *K $\rightarrow$ *ŋ}
\end{center}
The occurrence of this rule in the archaic strata of both Vedic and Old Tamil provides evidence for a shared phonological substrate that predates the great radiation from the Altai-Pamir hub.

